% Root literal audit and independent algebra audit by D: PASS.
\subsection{Projective endpoint coinvariants through seventeen arrows}

\begin{lemma}[Three degree-two classes at a five-dimensional quotient]
\label{C:three-by-four-five-outside-rank-three}
Let $P$ have rectangular-band minimal ideal of shape $3\times4$
or $4\times3$, with exactly five elements outside that ideal.
Suppose its rational homology is that of
$\mathbb CP^2\vee S^2\vee S^2$.  Then every product of two
rational degree-two classes on $BP$ is zero.  No finite-CW-type
or local-group-nullity hypothesis is required in this lemma.
\end{lemma}
\begin{proof}
After opposition the band--Tor modules have dimensions
$\dim V=2$, $\dim W=3$ over the five-dimensional contracted
algebra $A$, and their Tor dimensions are
\[
 \dim\operatorname{Tor}_{0,1,2}^A(W,V)=(3,0,1),\qquad
 \operatorname{Tor}_j^A(W,V)=0\quad(j\geq3).
\]
If $A$ is local, the dimension of its residue division algebra
divides both module dimensions, so that residue algebra is
$\mathbb Q$.  If $V$ has top dimension two, then $V$ is the
sum of two copies of the simple module; its degree-zero Tor
has even dimension, a contradiction.  Otherwise $V$ is cyclic.
Write $0\to I\to A\to V\to0$, where $I$ is nonzero.
Since degree-one Tor vanishes, tensoring gives an injection
$W\otimes_A I\to W$.  The following surjection
$W\to W\otimes_A V$ is an isomorphism, because both spaces
have dimension three.  Hence $W\otimes_A I=0$.
But this tensor product surjects onto
\[
 (W/WJ)\otimes_{\mathbb Q}(I/JI),
\]
which is nonzero by Nakayama's lemma.  This is impossible.

For nonlocal $A$, the Tor-persistence classification in
Proposition~\ref{B:five-nonlocal-exceptions} leaves the
zero-composite line and Nakayama algebras.  The complete
indecomposable calculation in Lemma~\ref{B:five-exceptional-targets}
applies before its particular module-dimension specialization.
For the line, it forces
\[
 V=S_1\oplus S_3,\qquad
 W=T_1\oplus T_3^{\oplus s}\oplus(e_1A)^{\oplus p}
                         \oplus(e_2A)^{\oplus q}.
\]
The equations are $s+2p+2q=2$ and $1+s+p=3$;
thus $W=T_1\oplus T_3^{\oplus2}$.
The actual three-point row augmentation is a sum of rational
characters, so Corollary~\ref{B:three-row-semisimple-cups}
makes all degree-two products zero.
For the Nakayama algebra the same table forces $V=L$ and
$W=T_2\oplus(e_1A)^{\oplus p}\oplus(e_2A)^{\oplus q}$.
Its dimension equation $3p+2q=2$ gives $p=0,q=1$,
and then degree-zero Tor has dimension one rather than three.
This excludes the last case and proves the lemma.
\end{proof}

\begin{theorem}[Projective endpoint coinvariants through seventeen arrows]
\label{C:semisimple-endpoints-seventeen}
Let $C$ have three linearly ordered objects, no backward arrows,
and nonempty forward hom-sets.  Suppose that the trivial left
module at the initial endpoint and the trivial right module at
the final endpoint are projective over their rational monoid
algebras.  If $C$ has at most seventeen arrows, then
$BC\not\simeq\mathbb CP^2$.
\end{theorem}
\begin{proof}
Suppose otherwise.  The actual endpoint-removal functors of
Corollary~\ref{C:finite-target-semisimple-endpoints} give a
triangle with trivial endpoints and sizes $(m;s,t,u)$.
Its compression has rational homology
$\mathbb CP^2\vee\bigvee_{u-1}S^2$, rational pairing of
rank one, and an integral class with nonzero primitive free
square.  Every integral degree-two class restricts to zero
on every finite group.  The budget is
\[
 m+s+t+u+2\leq17.
\]
The existing argument for sixteen arrows excludes $m\leq4$
and $\min(s,t)\leq2$ without any bound on the other sizes.
Thus take $m\geq5$, $3\leq s\leq t$, and $u\geq1$.
The entire remaining numerical scope is
\[
\begin{array}{c|c|c}
 (s,t)&\text{bounds}&\text{remaining treatment}\\\hline
 (3,3)&m+u\leq9&\text{all-rank zero products}\\
 (3,4)&m+u\leq8&u=3\text{ is the new case}\\
 (3,5)&m+u\leq7&\text{rank one or two}\\
 (3,6)&(m,u)=(5,1)&\text{five-outside module obstruction}\\
 (4,4)&m+u\leq7&\text{rank one or two}\\
 (4,5)&(m,u)=(5,1)&\text{five-outside module obstruction}.
\end{array}
\]
The first two rows use Lemma~\ref{C:small-band-finite-target-cups},
except for the newly proved rank-three case, which uses
Lemma~\ref{C:three-by-four-five-outside-rank-three}.
The rank-one $3\times5$ and $4\times4$ cases use
Propositions~\ref{B:three-by-five-finite-target-transfer} and
\ref{B:four-by-four-finite-target-transfer}; their compression
orders are at most twenty-one and twenty-two, respectively.
The two rank-two cases have compression orders twenty and
twenty-one, so Theorem~\ref{B:rank-two-rational-through22}
applies with the stated local nullity and rational homology.

Finally the proofs of
Lemmas~\ref{C:three-by-six-five-outside-CP2} and
\ref{C:four-by-five-five-outside-CP2} use only the rational
Tor dimensions $(1,0,1,0,\ldots)$ and zero rational squares,
not finite CW type.  In each, the local case has a rank-one
balanced form, and the two nonlocal algebras are excluded by
the displayed module equations or by the rank-one square
primitive.  Those proofs therefore apply to the present
finite-target compressions as well.
Every numerical case is excluded.  At no point is a finite
CW model inferred for an endpoint quotient or compression.
\end{proof}
